% Verze pro jednostranný tisk:
\documentclass[11pt,a4paper]{report}
\usepackage[top=25mm,bottom=25mm,right=25mm,left=30mm,head=12.5mm,foot=12.5mm]{geometry}
\let\openright=\clearpage

%% Definice různých užitečných maker (viz popis uvnitř souboru)
%%% Tento soubor obsahuje definice různých užitečných maker a prostředí %%%
%%% Další makra připisujte sem, ať nepřekáží v ostatních souborech.     %%%

\usepackage[a-2u]{pdfx}     % výsledné PDF bude ve standardu PDF/A-2u

%%% Nastavení pro použití samostatné bibliografické databáze.
\usepackage[
   backend=biber
%  ,style=iso-authoryear
  ,style=authoryear
  ,sortlocale=cs_CZ
  ,bibencoding=UTF8
  %,block=ragged
]{biblatex}
\let\cite\parencite
\bibliography{literatura}

%% Přepneme na českou sazbu, fonty Latin Modern a kódování češtiny
\usepackage[czech]{babel}
\usepackage{lmodern}
\usepackage[T1]{fontenc}
\usepackage{textcomp}
\usepackage[utf8]{inputenc}

%%% Další užitečné balíčky (jsou součástí běžných distribucí LaTeXu)
\usepackage{amsmath}        % rozšíření pro sazbu matematiky
\usepackage{amsfonts}       % matematické fonty
\usepackage{amsthm}         % sazba vět, definic apod.
\usepackage{bm}             % tučné symboly (příkaz \bm)
\usepackage{graphicx}       % vkládání obrázků
\usepackage{fancyvrb}       % vylepšené prostředí pro strojové písmo
\usepackage{fancyhdr}       % prostředí pohodlnější nastavení hlavy a paty stránek
\usepackage{icomma}         % inteligetní čárka v matematickém módu
\usepackage{dcolumn}        % lepší zarovnání sloupců v tabulkách
\usepackage{booktabs}       % lepší vodorovné linky v tabulkách
\makeatletter
\@ifpackageloaded{xcolor}{
   \@ifpackagewith{xcolor}{usenames}{}{\PassOptionsToPackage{usenames}{xcolor}}
  }{\usepackage[usenames]{xcolor}} % barevná sazba
\makeatother
\usepackage{multicol}       % práce s více sloupci na stránce
\usepackage{caption}
\usepackage{enumitem}
\setlist[itemize]{noitemsep, topsep=0pt, partopsep=0pt}
\setlist[enumerate]{noitemsep, topsep=0pt, partopsep=0pt}
\setlist[description]{noitemsep, topsep=0pt, partopsep=0pt}

\usepackage{tocloft}
\setlength\cftparskip{0pt}
\setlength\cftbeforechapskip{1.5ex}
\setlength\cftfigindent{0pt}
\setlength\cfttabindent{0pt}
\setlength\cftbeforeloftitleskip{0pt}
\setlength\cftbeforelottitleskip{0pt}
\setlength\cftbeforetoctitleskip{0pt}
\renewcommand{\cftlottitlefont}{\Huge\bfseries\sffamily}
\renewcommand{\cftloftitlefont}{\Huge\bfseries\sffamily}
\renewcommand{\cfttoctitlefont}{\Huge\bfseries\sffamily}

% vyznaceni odstavcu
\parindent=0pt
\parskip=11pt

% zakaz vdov a sirotku - jednoradkovych pocatku ci koncu odstavcu na prechodu mezi strankami
\clubpenalty=1000
\widowpenalty=1000
\displaywidowpenalty=1000

% nastaveni radkovani
\renewcommand{\baselinestretch}{1.20}

% nastaveni pro nadpisy - tucne a bezpatkove
\usepackage{sectsty}    
\allsectionsfont{\sffamily}

% nastavení hlavy a paty stránek
\renewcommand{\footrulewidth}{.5pt}
\fancypagestyle{plain}{%
\fancyhf{} % clear all header and footer fields
\fancyfoot[C]{\thepage} %% [RO,LE] vyvolalo warning, odstraněno
\renewcommand{\headrulewidth}{0pt}
\renewcommand{\footrulewidth}{0.5pt}}

% Tato makra přesvědčují mírně ošklivým trikem LaTeX, aby hlavičky kapitol
% sázel příčetněji a nevynechával nad nimi spoustu místa. Směle ignorujte.
\makeatletter
\def\@makechapterhead#1{
  {\parindent \z@ \raggedright \sffamily
   \Huge\bfseries \thechapter. #1
   \par\nobreak
   \vskip 20\p@
}}
\def\@makeschapterhead#1{
  {\parindent \z@ \raggedright \sffamily
   \Huge\bfseries #1
   \par\nobreak
   \vskip 20\p@
}}
\makeatother

% Trochu volnější nastavení dělení slov, než je default.
\lefthyphenmin=2
\righthyphenmin=2

% Zapne černé "slimáky" na koncích řádků, které přetekly, abychom si
% jich lépe všimli.
\overfullrule=1mm

%% Balíček hyperref, kterým jdou vyrábět klikací odkazy v PDF,
%% ale hlavně ho používáme k uložení metadat do PDF (včetně obsahu).
%% Většinu nastavítek přednastaví balíček pdfx.
\hypersetup{unicode}
\hypersetup{breaklinks=true}
\hypersetup{hidelinks}
\hypersetup{
    colorlinks = true,
    allcolors = black,
    urlcolor = blue
}

%%% Prostředí pro sazbu kódu, případně vstupu/výstupu počítačových
%%% programů. (Vyžaduje balíček fancyvrb -- fancy verbatim.)

\DefineVerbatimEnvironment{code}{Verbatim}{fontsize=\small, frame=single}

%%% Použití zkratek
\usepackage{acronym}
\usepackage{csquotes}

%%% Zabrání rozdělení paragrafů mezi stránkami
\widowpenalties 1 10000
\raggedbottom

\usepackage{float}
\usepackage{mathtools}
\DeclarePairedDelimiter{\ceil}{\lceil}{\rceil}

\graphicspath{ {./img/} }
\overfullrule=0pt


%%% Údaje o práci
% Název práce v jazyce práce (přesně podle zadání)
\def\NazevPrace{Modely logistické regrese v oblasti esportových dat}

% Typ práce
\def\TypPrace{BAKALÁŘSKÁ PRÁCE}

% Jméno autora
\def\AutorPrace{Michal Lauer}

% Rok odevzdání. měsíc (slovně)
\def\DatumOdevzdani{Duben 2022}

% Vedoucí práce: Jméno a příjmení s~tituly
\def\Vedouci{Ing. Zdeněk Šulc, Ph.D.}

% Studijní program a obor
\def\StudijniProgram{Aplikovaná informatika}
\def\StudijniObor{Aplikovaná informatika}

% Text čestného prohlášení pro MUŽE pro bakalářskou práci
\def\Prohlaseni{Prohlašuji, že jsem bakalářskou práci \textit{\NazevPrace} vypracoval samostatně za použití v práci uvedených pramenů a literatury.}

% Nepovinné poděkování (vedoucímu práce, konzultantovi, tomu, kdo
% zapůjčil software, literaturu apod.)
\def\Podekovani{%
Rád bych poděkoval panu doktoru Zdeňku Šulcovi, který mou bakalářskou práci podpořil a vedl, i přes odlišný studijní obor.
Dále děkuji autorům citovaných knih za poskytnutou příležitost se ve logistických modelech zlepšit. Bez nich by se práce psala velmi složitě.
}

% Abstrakt (doporučený rozsah cca 150-250 slov; nejedná se o zadání práce)
\def\Abstrakt{%
Cílem bakalářské práce je kvantitativně zanalyzovat esportové zápasy ze hry \ac{CSGO}, predikovat výhry
hráčů a týmů a zjistit, jaké prediktory jsou statisticky významné. Esport je moderní 
Použitý datový soubor je z internetového portálu
kaggle.com a obsahuje data od roku 2015 až do roku 2020. Predikce jsou založené na vícerozměrném logistickém regresním modelu.
Modely jsou sestavené jak pro hráče, tak pro týmy a to přes různé kategorie map. Modely jsou následně vyhodnoceny pomocí matice záměn a 
významnost koeficientů je testována pomocí Waldova testu. Modely jsou následně interpretovány v kontextu esportu a reálného života.

}

\def\AbstraktEN{%
The bachelor thesis aims to quantitatively analyze esport matches from the game \ac{CSGO}, predict the winnings of players and teams, and find out which predictors 
are statistically significant. The data file used is from the internet portal kaggle.com and contains data from 2015 to 2020. The thesis is divided into three parts. The first 
part introduces esport as such, summarizes its development, and defines important concepts and terms. The second part contains a description of the methods used for analysis 
and prediction. The predictions are based on a multidimensional logistic regression model. In the final practical part, the models are created for both players and teams across 
different categories of maps. The models are then evaluated using a confusion matrix, and the significance of the coefficients is tested using the Wald test. The models are 
then interpreted in the context of esports and real life.
}


% 3 až 5 klíčových slov (doporučeno)
\def\KlicovaSlova{model, logistická regrese, predikce, esport}
\def\KlicovaSlovaEN{model, logistic regression, prediction, esport}

%% Titulní strana a různé povinné informační strany
\begin{document}
\pagenumbering{gobble}
%%% Titulní strana práce a další povinné informační strany

%%% Titulní strana práce

\pagestyle{empty}
\hypersetup{pageanchor=false}

\begin{center}
\Huge\sffamily
Vysoká škola ekonomická v Praze\\
Fakulta informatiky a statistiky

\vspace{\stretch{1}}

\includegraphics[width=.3\textwidth]{img/logo-FIS}

\vspace{\stretch{2}}

\bfseries\NazevPrace

\vspace{8mm}
\mdseries\TypPrace

\vspace{8mm}
\large
\begin{tabular}{rl}
Studijní program: & \StudijniProgram \\
\noalign{\vspace{2mm}}
Studijní obor: & \StudijniObor \\
\end{tabular}

\vspace{\stretch{8}}

\begin{tabular}{rl}
Autor: & \AutorPrace \\
\noalign{\vspace{2mm}}
Vedoucí práce: & \Vedouci \\
\end{tabular}

\vspace{8mm}
Praha, \DatumOdevzdani
\end{center}


\openright

%%% Strana s čestným prohlášením k bakalářské práci

\hypersetup{pageanchor=true}
\pagestyle{plain}
\cleardoublepage
\vspace*{\fill}
\section*{Prohlášení}
\noindent
\Prohlaseni

\vspace{2cm}
\noindent
V Praze dne DD. Dubna 2021
\hfill%
\begin{minipage}[t]{.5\textwidth}%
\begin{center}
\dotfill\\
Podpis studenta
\end{center}
\end{minipage}
\vspace{1cm}

%%% Poděkování
\openright
\vspace*{\fill}
\section*{Poděkování}
\noindent
\Podekovani
\vspace{1cm}


%%% Povinná informační strana bakalářské práce
\openright
\section*{Abstrakt}
\noindent
\Abstrakt
\subsection*{Klíčová slova}
\noindent
\KlicovaSlova

\bigskip\bigskip\bigskip
\section*{Abstract}
\noindent
\AbstraktEN
\subsection*{Keywords}
\noindent
\KlicovaSlovaEN

\openright


%%% Strana s automaticky generovaným obsahem bakalářské práce
\setcounter{tocdepth}{2}
\tableofcontents

\pagestyle{fancy}
\pagenumbering{arabic}
%%% Jednotlivé kapitoly práce jsou pro přehlednost uloženy v samostatných souborech
\chapter{Úvod}
Esport je označení pro elektronický sport. Obsahuje všechny důležité oblasti jako klasický sport (např. turnaje, trénování, investice, stadiony, či sázení)
s tím rozdílem, že se hraje na nějakém zařízení (počítač, konzole, mobil).
Je to jedno z nejrychleji rostoucích odvětví v dnešní době. V roce 2021 se tržní hodnota esportu pohybovala kolem jedné miliardy dolarů - skoro
50\% nárůst oproti roku 2020. Lze předpovídat, že v roce 2024 esport překročí hodnotu 1,5 miliardy dolarů \cite{Gough2021}.
Dalo by se spekulovat, že za takový velký nárůst je zodpovědná aktuální pandemie. Většina populace je nucena zůstat doma. Toto otevřelo dveře
se s esportem přirozeně seznámit a nějakým způsobem se ho účastnit (online divák, soutěžící, organizátor, fanoušek...). 
Hrají se různé kategorie her jako např. střílečky, \ac{MOBA}\footnote{tzn. MOBA, kde hráči hrají v jedné online aréně proti sobě},
karetní hry, \ac{FPS} či \ac{BR}.

Práce se zaměřuje primární na esportový titul \acf{CSGO}. Je to jeden z nejdéle hraných esportových titulů, boří mnohé divácké
rekordy\footnote{\scriptsize \url{https://www.invenglobal.com/articles/15619/csgo-major-breaks-viewership-records-overtakes-the-international}}
a je aktuálně nejhranějším \ac{FPS} esport titulem. \ac{CSGO} vyniká nejen detailní herní mechanikou, ale i bohatou a zajímavou historií. Hra
je unikátní i tím, že obsahuje mnoho různých módu\footnote{rozšíření, jak hru hrát. Každý mód má svá vlastní pravidla, mapy, či herní fanoušky}
a hráč může strávit mnoho hodin pouze objevováním komunitních serverů, hraním klasických zápasů či trénováním na offline mapách.

Finální cíl práce je vytvořit logistický regresní model, který předpovídá výsledek zápasů. Pro tvorbu kvalitního modelu bude kritické zvolit vhodné prediktory.
Pro predikci jsou použity prediktory, které se nacházejí ve dvou samostatných datových souborech\footnote{\url{https://www.kaggle.com/mateusdmachado/csgo-professional-matches}} 
které podávají informace jak už o zápase (např. datum, 
výsledek zápasu, výsledek jednotlivých map, typ zápasu), hráčích (např. statistiky za zápas, statistiky za mapy, statistiky za team), tak o 
vývoji celého zápasu (především ekonomika týmu). V práci bude tedy vytvořeno více specializovaných modelů pro každý vybraný týmů a následně
je pro každý tým vybrán nejlepší model. Výsledné modely jsou v závěru mezi sebou porovnány.

Logistický model je preferován kvůli své lehké interpretaci a 
dobré aplikaci v reálném životě. Výsledky, statistiky a pravděpodobnosti mohou být použity např. v sázkových kancelářích, kdy se výsledky modelu
dají využít na nejrůznější sázky a lze předpovídat,
kdo vyhraje zápas, kdo vyhraje mapu, jaký hráč bude mít nejlepší statistiky, či zda si hráč koupí určitou zbraň.

Práce je tedy rozdělená do tří částí. V první části je kladen důraz na esport, jeho vývoj, a na esportový titul \ac{CSGO}. Jsou zde také představená pravidla, podle kterých se
hra hraje. V druhé části jsou popsány popisné a statistické metody. Jsou zde definované grafické nástroje pro popis datasetu, logistický regresní model, a evaluační nástroje
pro model. Třetí část se zaměřuje na praktickou tvorbu modelů, jejich interpretaci, a vzájemné porovnání.

\chapter{Představení esportu}
\section{Historie esportu}
\section{Zasazení do dnešní doby}
\section{Představení vybraného esportového titulu}
\section{Propojení bakalářské práce a esportového titulu}
\chapter{Teoretická část}
V následující části jsou popsány jak teoretické metody pro vizualizaci dat, tak i tvar, forma a vyhodnocení logistického regresního modelu. 
Ke každé části, která se věnuje popisu dat pomocí nějakého grafu, je přidána praktická ukázka s popisem a praktickým vysvětlením. U regresního 
modelu je následně dán důraz na teorii, jelikož model je následně prakticky použit v celé třetí části. Důležité je si uvědomit, že nástrojů pro
vizualizaci dat je mnoho a grafů či souhrnných statistik je nespočet. V této práci se ovšem věnujeme vizualizacím, které jsou pro regresní model
vhodné.
Testovací citace: \cite{Hebak2015}, \cite{Kleinbaum2010}

\section{Vizualizace dat}
\subsection{Bodový graf}
Bodový graf slouží pro zobrazení vztahu dvou numerických proměnných. Z pravidla se vysvětlovaná proměnná dává na osu Y,
zatímco proměnná vysvětlující se nachází na ose X. Graf \ref{fig:bodovy_graf_mtcars} zobrazuje negativní vztah mezi
váhou vozidla a mílemi ujetými za galon.

\begin{figure}[H]
    \centering
    \includegraphics{../../analyza/plots/bodovy_graf_mtcars.png}
    \caption{Bodový graf váhy a míly za galon z data setu mtcars} 
    \label{fig:bodovy_graf_mtcars}
\end{figure}

\subsection{Histogram}

Histogram je speciální typ sloupcového grafu. Je specifický tím, že zobrazuje počet jednotlivých pozorování ve skupině. Skupiny je nutné vytvořit
v případě, že je zobrazovaná proměnná spojitá. Zároveň histogram oproti klasickému sloupcovému grafu, který zobrazuje kategorie, nemá mezi sloupci žádné mezery.
Na grafu \ref{fig:sloupcovy_graf_mtcars} je příklad sloupcového grafu. Na ose X se nachází kategorická vysvětlovací proměnná. Vysvětlovaná proměnná Y je průměrná hrubá
koňská síla.

\begin{figure}[H]
    \centering
    \includegraphics{../../analyza/plots/sloupcovy_graf_mtcars.png}
    \caption{Sloupcový graf z data setu mtcars} 
    \label{fig:sloupcovy_graf_mtcars}
\end{figure}

Pro histogram je třeba data sloučit do skupin \textit{(bins)} o určité šířce. Správný výběr počtu skupin je kritický, jelikož může velmi
silně ovlivnit interpretaci dat. Pokud se vybere moc malý počet skupin, data se seskupí a může se ztratit důležitý vztah. Pokud se ovšem
vybere moc velký počet skupin, v datech bude obtížné najít nějaký obecný vztah či trend.
Tento efekt je znázorněn na grafu \ref{fig:histogram_porovnani_mtcars}.

\begin{figure}[H]
    \centering
    \includegraphics{../../analyza/plots/histogram_porovnani_mtcars.png}
    \caption{Porovnání histogramů s různým počtem skupin} 
    \label{fig:histogram_porovnani_mtcars}
\end{figure}

Pro vhodný počet skupin existuje mnoho způsobů. Nejznámější je takzvané Sturgesovo pravidlo (\ref{eq:sturgesovo_pravidlo}), které se spočítá následujícím vztahem:

\begin{equation}
    \label{eq:sturgesovo_pravidlo}
    k \text{ } \dot{\mathbf{=}} \text{ } 1 + 3,22 * log_{10}(n)
\end{equation}

kde $k$ je výsledný zaokrouhlený počet skupin a $n$ je počet pozorování. Druhý parametr, který je pro tvorbu histogramu potřeba, je šířka skupiny.
Ta by měla být ideálně stejná pro všechny skupiny. Pokud tomu tak není, histogram může být zavádějící a čtenář mu nemusí plně rozumět.
Pro vypočtení počtu skupin má šířka skupiny následující tvar:

\begin{equation}
    \label{eq:sirka_histogramu}
    w = \ceil[\bigg]{ \frac{max(x) - min(x)}{k} }
\end{equation}

kde $x$ je zobrazovaná proměnná, $k$ je počet skupin a $w$ je výsledná šířka intervalu, která je zaokrouhlena nahoru. To je z toho důvodu, aby se
každé pozorování promítlo právě do jedné kategorie. Pokud na stejný dataset, jako na grafu \ref{fig:histogram_mtcars_sturges}, použije sturgesovo
(rovnice \ref{eq:sturgesovo_pravidlo}) pravidlo, graf vypadá následovně:

\begin{figure}[H]
    \centering
    \includegraphics{../../analyza/plots/histogram_mtcars_sturges.png}
    \caption{Histogram s počtem skupin dle Sturgesova pravidla} 
    \label{fig:histogram_mtcars_sturges}
\end{figure}

Histogram lze samozřejmě rozepsat. Lze vytvořit tzn. tabulku četností (\ref{tab:tabulka_cetnosti_sturges}),
která mimo jiné obsahuje spodní hranici intervalu, horní hranici intervalu a počet pozorování. Důležité je,
aby přechody mezi intervaly byli jasné.

\begin{table}[H]

        \centering
    \begin{tabular}[t]{c|c|c}
        \hline
        Spodní hranice intervalu & Horní hranice intervalu & Počet pozorování\\
        \hline
        7.05 & 11.75 & 2\\
        \hline
        11.76 & 16.45 & 9\\
        \hline
        16.46 & 21.15 & 9\\
        \hline
        21.16 & 25.85 & 6\\
        \hline
        25.86 & 30.55 & 4\\
        \hline
        30.56 & 35.25 & 2\\
        \hline
    \end{tabular}
    \caption{\label{tab:tabulka_cetnosti_sturges}Tabulka četnostní}
\end{table}

\subsection{Boxplot}
\subsubsection{Five-number summary}
Five-number summary je číselná tabulka, která pomocí pěti různých čísel shrnuje seřazenou číselnou řadu. Základní statistický nástroj pro
vytvoření takové tabulky jsou percentily. Hodnota $P$-tého percentilu označuje číslo, které rozděluje seřazenou číselnou řadu na dva intervaly. 
První interval obsahuje $P*100\%$ číselné řady a druhý analogicky $(1-P)*100\%$. Různé hodnoty percentilů mohou mít specifičtější pojmenování a značí se $Q_P$.
Percentil $P = 0.5$ se označuje jako medián a rozděluje seřazenou číselnou řadu na polovinu. Percentily, kde $P = 0.25$ nebo $P = 0.75$, se označují
jako kvartily a značí se $Q_{1}$ a $Q_{3}$. Oba tyto typy kvartilů jsou použité při tvorbě Five-number summary tabulky. Jako příklad je
uvedena tabulka \ref{tab:five-number_summary}

\begin{table}[H]

    \centering
    \begin{tabular}[t]{c|c|c|c|c}
        \hline
        $Q_{0} (Q_0)$ & $Q_{0.25} (Q_1) $ & $Q_{0.50}$ & $Q_{0.75} (Q_3)$ & $Q_{1.00}$\\
        \hline
        1.513 & 2.58125 & 3.325 & 3.61 & 5.424\\
        \hline
    \end{tabular}
    \caption{\label{tab:five-number_summary}Five-number summary tabulka hmotnosti vozidla (lb/1000)}
\end{table}

$Q_{0}$ a $Q_{1.00}$ označují minimum a maximum číselné řady. Kvartily $Q_{1}$ a $Q_{3}$ jsou čísla, která rozděluji časovou řadu na na čtvrtiny. V prvním
případě, tedy $Q_1 = Q_{0.25}$, je 25\% čísel menší než 1.513 a 75\% dat větší. Pro kvantil $Q_3 = Q_{0.75}$ je 75\% čísel menších než 3.61 a 25\% větších. $Q_{0.50}$ označuje medián.

\subsubsection{Boxplot}
Boxplot je grafické zobrazení a rozšíření Five-number summary tabulky. Krom grafické zobrazení pěti percentilů obvykle ukazuje statistiky jako průměr a extrémní hodnoty.
V boxplotu se také nachází obdélník, který ukazuje interkvartilní rozsah \textit{(IQR)}, tedy prostředních 50\% dat. V obdélníku se také nachází černá čára, která značí medián.
Z prostředního obdélníku vedou oběma směry čáry, které značí teoretické minimum a maximum (ne hodnoty $Q_0$ a $Q_1$). Obě dve hranice se počítají následovně

\begin{align}
    min &= Q_1 - 1.5 * IQR \\
    max &= Q_3 + 1.5 * IQR
\end{align}

Hodnoty, které spadají do intervalu $<Q_1 - 1.5 * IQR ; Q_1>$ a $<Q_3, Q_3 + 1.5 * IQR>$ se nazývají jako odlehlé. Hodnoty které leží mimo tento vztah,
tedy hodnoty menší než $Q_1 - 1.5 * IQR$ či větší než $Q_3 + 1.5 * IQR$, mají speciální název a označují se jako hodnoty extrémní.
V boxplotu jsou z pravidla znázorněni nějakým speciálním znakem, např. kolečkem.
Díky grafickém zobrazení lze lehce porovnávat distribuci jedné vysvětlované proměnné vůči několika kategoriím.

\begin{figure}[H]
    \centering
    \includegraphics{../../analyza/plots/boxplot_mtcars.png}
    \caption{Boxplot váhy auta pro různé počty válců} 
    \label{fig:boxplot_mtcars}
\end{figure}

Trojúhelník v grafu \ref{fig:boxplot_mtcars} značí průměr pro danou kategorii. Černé tečky v kategorii osmi válců značí extrémní hodnoty, t. j. hodnoty
větší než $<Q_3, Q_3 + 1.5 * IQR>$ jsou extrémní.

\section{Logistický regresní model}
... Dopsat

\subsection{Jednoduchý logistický regresní model}
... Dopsat

\subsection{Vícerozměrný logistický regresní model}
... Dopsat

\subsection{Optimalizace parametrů}
... Dopsat

\subsubsection{Metoda nejmenších čtverců}
... Dopsat

\subsubsection{Metoda největší pravděpodobnosti}
... Dopsat

\subsection{Vyhodnocení modelu}
... Dopsat

\chapter{Praktická část}

{\color{red}
V této části se nachází popis dat, transformace dat, a tvorba logistického regresního modelu. Nejprve jsou představené datové soubory,
se kterými se pracuje. Následně jsou použité grafy, které jsou představené v sekci \ref{sec:vizualizace_dat}. Pomocí vizualizace dat lze představit proměnné, které
by mohli do logistického modelu vstupovat. Také lze díky grafům získat intuici, jak datový soubor vypadá a jaké výsledky je možné očekávat.
Poté jsou vytvořené logistické regresní modely, jejich výstup je interpretován a různé modelu jsou mezi sebou porovnány. V závěru každé podsekce
se nachází zamyšlení, jak by daný model mohl být vylepšen a jaké je jeho použití v reálném životě.
}

{\color{red}
\section{Datové soubory}
Dataset\footnote{https://www.kaggle.com/datasets/mateusdmachado/csgo-professional-matches} obsahuje čtyři souboru, které popisují zápasy ve hře
\ac{CSGO}. K potřebám této bakalářské práce budou použity pouze soubory \textit{players.csv} a \textit{results.csv}. Ostatní dva soubory obsahují
informace, které nezle k predikci v případě této bakalářské práce použít.
}

{\color{red}
\subsection{soubor players.csv}
Soubory \textit{players.csv} obsahuje statistiky jednotlivých hráčů v daném zápase. Původní dataset obsahuje 101 proměnných a 379 680 pozorování.
V původním datovém souboru se jeden řádek (záznam) rovná statistikám jednoho hráče za celý zápas. Pro potřeby bakalářské práce je však nutné získat statistiky hráčů
za jednotlivé mapy (jeden zápas může být hrán až na třech mapách). Proto je původní datový soubor transformován do podoby, kde se jeden záznam rovná statistikám
právě jednoho hráče na právě jedné mapě, a to bez ohledu, kolik map se v daném zápase hrálo. Jinak řečeno, transformovaný datový soubor nebere v potaz, zda
se daná mapa hrála jako první, druhá, či třetí.
Transformovaný dataset má 10 proměnných a 643 620 záznamů. Příklad jednotlivých pozorování v transformovaném
datasetu je v přiložené tabulce \ref{tab:players_csv_transformovano}.
}
\newpage
Transformovaný dataset má 10 proměnných, které unikátně identifikují statistiky každého hráče na určité mapě v jednom zápase. Interpretace je následující:
\begin{itemize}
    \item \textbf{match\_id} --- identifikátor zápasu
    \item \textbf{player\_id} --- identifikátor hráče
    \item \textbf{team} --- jméno týmu
    \item \textbf{map} --- název hrané mapy
    \item \textbf{kills} --- počet zabití hráče v zápase na dané mapě
    \item \textbf{assists} --- počet asistencí hráče v zápase na dané mapě
    \item \textbf{deaths} --- počet smrtí hráče v zápase na dané mapě
    \item \textbf{hs} --- procento zabití, které lze označit jako headshot\footnote{hráč zabil nepřítele střelou do hlavy}
    \item \textbf{fkdiff} --- rozdíl, kolikrát hráč zabil jako první nepřítele versus kolikrát byl jako první zabit
    \item \textbf{rating} --- shrnutí více statistik za jeden zápas do jednoho ukazatele výkonu\footnote{\url{https://www.hltv.org/news/20695/introducing-rating-20}}
\end{itemize}

{\color{red}
\subsection{soubor results.csv}
Druhý datový soubor, který je pro analýzu použit, obsahuje výsledky daných zápasů. Dataset se původně skládá z 45 773 řádků a 19 proměnných. Dataset obsahuje na rozdíl
od datového souboru \textit{players.csv} jeden chybný záznam. Dle něho hrál tým sám proti sobě, což nedává smysl. Jelikož je zápas na webovém portálu zadán
správně, nejspíše se jedná o neznámou chybu, která nastala při exportu dat z webového portálu.
}
Po transformacích vznikne tabulka o 8 proměnných a 91 436 řádcích. Každé pozorování identifikuje výsledek jednoho týmu v jednom zápasu
na jedné mapě. Příklad je zobrazen v přiložené tabulce \ref{tab:results_csv_transformovano}. Jednotlivé proměnné lze interpretovat následovně:
\begin{itemize}
    \item \textbf{date} --- datum, kdy se hrál zápas
    \item \textbf{match\_id} --- identifikátor zápasu
    \item \textbf{team} --- jméno týmu
    \item \textbf{map} --- název hrané mapy
    \item \textbf{map\_winner} --- binární značení, zda tým vyhrál (1) či prohrál (0)
    \item \textbf{starting\_ct} --- binární značení, zda tým začal zápas na straně Counter-Teroristů (1) či Teroristů (0)
    \item \textbf{team\_rank} --- rank týmu v okamžik, kdy se zápas hrál\footnote{\url{https://www.hltv.org/news/16061/introducing-csgo-team-ranking}}
    \item \textbf{run\_mean\_3\_months} --- klouzavý průměr týmu za poslední tři měsíce
\end{itemize}

{\color{red}
\subsection{Omezení datasetu}
Dataset obsahuje pozorování o zápasech a statistikách od konce roku 2015 do začátku roku 2020. Díky velkému počtu záznamů je možné chybné či neúplné záznamy
smazat. To může nastat např. když tým má méně než 5 hráčů nebo když zápas nemá v každém týmu právě 5 hráčů. Tým může mít méně než 5 hráčů z toho důvodu, že je např.
amatérsky\footnote{neprofesionální, tím pádem nemusí mít všichni hráči na webovém portálu založená profil}. Pokud má hráčů více, vyberou se hráči dle délky svého působení v týmu.
Tým může mít v zápasu více než pět hráčů, pokud použijí náhradníka. Také jsou smazané záznamy, které neobsahují všechny potřebné statistiky. To může nastat u velmi historických
zápasů, kde výpočet daných statistik ještě nebyl možný.
}

\newpage
\section{Cíl analýzy}
Cílem analýzy je vytvořit logistické regresní modely, které předpovídají pravděpodobnost výhry daného hráče či týmu. Pro předpověď se použijí statistiky hráčů
za každou mapu a výsledné týmové statistiky za každou mapu. Statistiky hráčů, které jsou označovaný za prediktory jsou: \textit{kills}, \textit{assists}, \textit{deaths},
\textit{hs}, \textit{fkdiff} a \textit{rating}. Za týmové statistiky jsou pak považovány proměnné \textit{starting\_ct} a \textit{run\_mean\_3\_months}. Lze očekávat, že
prediktory \textit{kills}, \textit{assists}, \textit{hs}, \textit{rating} a \textit{fkdiff} budou šanci na výhru zvyšovat. Naopak prediktor \textit{deaths} by měl
výslednou šanci na výhru snížit. 

{\color{green}
Nehodila by se tato část spíše někde v úvodu bakalářské práce?
}
První typ logistického modelu bude predikovat výhru zápasu na určité mapě podle individuálních statistik hráče a týmových prediktorů. Nebere se v potaz
výkon spoluhráčů ani průměrný rank týmu za poslední tři měsíce. Cílem modelu je identifikovat, která statistika hráčů má na určité mapě největší vliv na výhru.

Druhý model se bude soustředit na predikci výhry týmu. K tom využije agregované statistiky hráčů, kteří jsou součástí týmu. Může se stát, že za tým historicky hrálo více hráčů.
V případě, že tým vymění jednoho hráče za druhého, bude za tým historicky hrát 6 hráčů. V takovém případě model počítá s hráči, kteří za tým hrají nejdelší dobu. V případě schody
se vezme hráč, který hrál v týmu jako poslední. Agregace statistik je spočítaná jako aritmetický průměr přes všechny zápasy, kde daná pětice za tým hraje. Dále jsou v modelu
prediktory \textit{starting\_ct} a \textit{run\_mean\_3\_months}. Cíl druhého modelu je zjistit, jaké prediktory jsou pro předpověď výhry týmu významné. Model je vytvořen
jak pro celý datový soubor, tak pro jeden vybraný tým.

Poslední třetí logistický model je vytvořen pro jeden specifický tým, který bude předpovídat pravděpodobnost výhry proti historicky nejlepším třiceti týmům. Cílem tohoto modelu
je zjistit, jak se liší významnost prediktorů, pokud tým hraje pouze proti těm nejlépe hodnocením týmům. 

\section{Průzkumová analýza dat}
Průzkumová analýza slouží k vizualizaci prediktorů, hledání různých vztahů a rozdělení proměnných. Díky průzkumu lze určit, které proměnné není vhodné použít pro tvorbu
logistického regresního modelu, např. kvůli problému multikolinearity.

\newpage
\subsection{Korelační matice}
Pro logistickou regresi je důležité, aby prediktory nebyli lineárně závislé. Kombinaci korelací mezi kvantitativními prediktory lze zjistit z korelační matice.

\begin{figure}[H]
    \centering
    \includegraphics{../obrazky/prediktory_corr_matice.png}
    \caption{Korelační matice} 
    \label{fig:korelacni_matice}
\end{figure}

Z korelační matice \ref{fig:korelacni_matice} lze vyčíst, korelace se blíží nule mezi průměrem týmu za poslední tři měsíce a agregovanými statistikami hráčů. Zároveň je
vidět středně silná korelace mezi prediktorem \textit{rating} a statistikami \textit{fkdiff}, \textit{deaths} a \textit{kills}. Z tohoto důvodu je
prediktor \textit{rating} z tvorby modelů vyloučen.

\newpage
\subsection{Histogram kvantitativních prediktorů}
Histogram kvantitativních prediktorů umožní zobrazit rozdělení prediktorů.

\begin{figure}[H]
    \centering
    \includegraphics{../obrazky/histogram_prediktoru.png}
    \caption{Histogram prediktorů} 
    \label{fig:histogram_prediktoru}
\end{figure}

Histogram prediktorů \textit{rating}, \textit{hs}, \textit{kills} a \textit{deaths} má normální rozdělení a v proměnné se nenachází mnoho extrémních hodnot.
Prediktor \textit{fkdiff} má bimodální\footnote{rozdělení má dva vrcholy} rozdělení. Prediktor \textit{assists} je sklonění doprava, což značí velké množství 
odlehlých či extrémních hodnot.

\newpage
\subsection{Sloupcový graf výher přes počáteční stranu}
Strana, na které tým začíná, může ovlivnit výsledek zápasu. Ovlivnění navíc může být rozdílné podle toho, na jaké mapě se zápas odehrává.

\begin{figure}[H]
    \centering
    \includegraphics{../obrazky/sloupce_podle_strany.png}
    \caption{Procento vyhraných zápasů na dané mapě za stranu Counter-Terroristů} 
    \label{fig:sloupcovy_graf_strany}
\end{figure}

Počáteční strana má největší vliv na mapě Cobblestone, kde necelých 52 procent týmů, co začalo na straně Counter-Terroristů, mapu vyhrálo. Sloupcový graf naopak naznačuje, že 
je pro týmy nevýhodné začínat na mapě Overpass za stranu Counter-Terroristů. Historicky necelých čtyřicet-devět procent týmu, co na mapě začalo za Counter-Terroristy, mapu prohrálo.

\section{Modely pro výkon hráče}
První vytvořený logistický regresní model se je založen pouze na výkonu hráče. Model je vytvořen na datech z obou datových souborů \textit{players.csv} a \textit{results.csv}.
Prediktory se týkají pouze výkonu jednotlivého hráče, model tedy pro předpověď výhry hráče nepoužívá statistiky spoluhráčů. Pro porovnání jsou vybrané mapy Cobblestone a Vertigo.
Mapa Mirage je jednou z nejvíce tradičních map a mapa Vertigo je naopak nejnovější přídavek do hry.

{\color{green}
Výběr prediktorů již nepopisuji, jelikož jsou vysvětlené v sekci Cíl analýzy
}

\subsection{Model pro mapu Mirage}
\input{kod/mirage_pro_hrace.tex}
Pro model jsou významné všechny prediktory. S rostoucím počtem zabití, asistencí a statistikou fkdiff má hráč větší šanci mapu vyhrát. Naopak s rostoucím počtem smrtí a 
statistikou hs, šance na výhru klesá. Mapa je pro stranu Counter-Terroristů nevýhodná, a pokud hráč začne mapu na této straně, jeho šance na výhru klesá. Největší vliv
má na pravděpodobnost výhry počet úmrtí hráče. Přepis funkce pak vypadá následovně
\begin{align}
    \begin{split}
        &P(1 | X_{kills}, X_{assists}, X_{deaths}, X_{hs}, X_{fkdiff}, X_{starting_ct}) = \frac{1}{1 + e^{-z}} \\
        &z = 1.7830 + 0.1820*X_{kills} + 0.2986*X_{assists} - 0.3504*X_{deaths} + \\
        &+ 0.1447*X_{hs} + 0.0064*X_{fkdiff} + 0.0610*X_{starting_ct}
    \end{split}
\end{align}



\subsection{Model pro mapu Vertigo}
\input{kod/vertigo_pro_hrace.tex}

Pro hráče jsou na mapě Vertigo významné pouze statistiky zabití, asistence a úmrtí. Ostatní prediktory hs, fkdiff a starting\_ct jsou pro model nevýznamné. Výsledek lze 
interpretovat tak, že pro hráče není významné, na jaké straně mapu začne hrát (prediktor  \textit{tarting\_ct}), ani jak přesně střílí (prediktor \textit{HS}).
Zde má na výhru hráče největší negativní vliv počet úmrtí. Po vyřazení nevýznamných prediktorů má model následující významné parametry

\input{kod/vertigo_pro_hrace_opt.tex}

Přepis funkce je pak následující
\begin{align}
    \begin{split}
        &P(1 | X_{kills}, X_{assists}, X_{deaths}) = \frac{1}{1 + e^{-z}} \\
        &z = 1.8792 + 0.1829*X_{kills} + 0.2985*X_{assists} - 0.3512*X_{deaths}
    \end{split}
\end{align}

\subsection{Matice záměn}
Pro porovnání modelů lze použít matice záměn a výkoností parametry, které z ní lze získat.

\begin{tabular}[t]{l|c|c}
    \hline
      & Predikovaná Výhra & Predikovaná prohra\\
    \hline
    Výhra & 35159 & 9672\\
    \hline
    Prohra & 8026 & 33513\\
    \hline
    \end{tabular}

\chapter{Závěr}
Esport neustále roste a proniká do odvětví jako je sázení, analytika či podnikání. Jedná se o nový trend, který nejvíce zajímá  mladou populaci a hráče her.
Propojením esportu a sázkových kanceláří vzniká potřeba pro kvalitní a přesné modely, díky kterým lze výsledky zápasů předpovídat, a stanovovat tak
vhodné kurzy. Na profesionální scéně lze modely využít k identifikaci významných prediktorů pro výhru v zápase. Hráči a týmy se pak díky datům a modelům mohou soustředit
pouze na charakteristiky, které jejich šanci na výhru zvyšují.

Vytvořené modely v bakalářské práci byly sestavené pro individuální hráče, pro dva referenční týmy i pro všechny týmy v datovém souboru. Z modelů
je patrné, že nejdůležitější prediktory se týkají počtu zabití nepřátel a počtu úmrtí hráče. Ty hrají významnou roli jak v modelech pro hráče, tak i pro tým. Rozdíl ve
významnosti prediktorů se nachází mezi profesionálními a amatérskými týmy. Z modelů lze odvodit důraz na individuální charakteristiky hráčů v profesionálních týmech. V
týmech amatérských lze pak pozorovat odvodit význmanost charakteristik \textit{mean\_assists}, \textit{mean\_hs}, \textit{mean\_fkdiff}, \textit{team\_rank} a významnost
interakce mezi mapou (\textit{map}) a počáteční stranou (\textit{starting\_ct})

Model by bylo možné rozšířit o více podrobných a specifických charakteristik hráčů. Modely navíc nepočítají s tím, že se hráči v týmu mohou obměňovat a že každý hráč má jiné
charakteristiky a je součástí týmu po rozdílnou dobu. Aktuální modely považují za známé, na jaké mapě se zápas hraje a jakou rolí tým začíná. V praxi je však tento fakt známý
až velmi krátce před začátkem zápasu. K modelům by bylo proto vhodné přidat predikci mapy, která se bude mezi dvěma týmy hrát.

%%% Seznam použité literatury
%% Toto platí v případě použití samostatné bibliografické databáze
\printbibliography[title={Seznam použité literatury}, heading={bibintoc}, keyword={literatura}]

\printbibliography[title={Seznam elektronických zdrojů}, heading={bibintoc}, keyword={ezdroj}]


%%% Obrázky v bakalářské práci
\openright
\phantomsection
\addcontentsline{toc}{chapter}{\listfigurename}
\listoffigures

%%% Tabulky v bakalářské práci (opět nemusí být nutné uvádět)
\clearpage
\phantomsection
\addcontentsline{toc}{chapter}{\listtablename}
\listoftables

%%% Použité zkratky v bakalářské práci (opět nemusí být nutné uvádět)
\chapter*{Seznam použitých zkratek}
\addcontentsline{toc}{chapter}{Seznam použitých zkratek}

\begin{acronym}
    %% esport tituly
    \acro{CSGO}{Coutner-Strike: Global Offensive}
    %% esport žánry
    \acro{BR}{Battle Royale}
    \acro{MOBA}{Multiplayer Online Battle Arena}
    \acro{FPS}{First-Person Shooter}
    %% ostatní
    \acro{TGNS}{Twin Galaxies National Scoreboard}
\end{acronym}


%%% Přílohy k bakalářské práci
\part{Přílohy}
\appendix
\chapter{Datové soubory}
\section{Transformovaný datový soubor players.csv}
\begin{table}[H]

\caption{\label{tab:players_csv_transformovano}Záznam z transformovaného datového souboru players.csv}
\centering
\resizebox{\linewidth}{!}{
\begin{tabular}[t]{c|c|c|c|c|c|c|c|c|c}
\hline
match\_id & player\_id & team & map & kills & assists & deaths & hs & fkdiff & rating\\
\hline
2339385 & 8738 & Liquid & Overpass & 15 & 3 & 12 & 0.6 & 3 & 1.32\\
\hline
\end{tabular}}
\end{table}


\section{Transformovaný datový soubor results.csv}
\begin{table}[H]

\caption{\label{tab:results_csv_transformovano}Příklad záznamu z transformovaného datového souboru results.csv}
\centering
\resizebox{\linewidth}{!}{
\begin{tabular}[t]{c|c|c|c|c|c|c|c}
\hline
date & match\_id & team & map & map\_winner & starting\_ct & team\_rank & run\_mean\_3\_months\\
\hline
2019-11-07 & 2337454 & 100 Thieves & Nuke & 0 & 1 & 8 & 8\\
\hline
\end{tabular}}
\end{table}


\section{Spojený datový soubor}
\input{tabulky/data_spojena.tex}

\section{Agregovaný datový soubor}
\input{tabulky/data_agregovana.tex}
\chapter{Modely, matice záměn a statistiky pro individuální hráče} \label{apx:02}
\section{Mapa Cache}
\input{modely/Cache_model.tex}
\input{matice/Cache_matice.tex}
\input{statistiky/Cache_stats.tex}
\newpage

\section{Mapa Cobblestone}
\input{modely/Cobblestone_model.tex}
\input{matice/Cobblestone_matice.tex}
\input{statistiky/Cobblestone_stats.tex}
\newpage

\section{Mapa Dust2}
\input{modely/Dust2_model.tex}
\input{matice/Dust2_matice.tex}
\input{statistiky/Dust2_stats.tex}
\newpage

\section{Mapa Inferno}
\input{modely/Inferno_model.tex}
\input{matice/Inferno_matice.tex}
\input{statistiky/Inferno_stats.tex}
\newpage

\section{Mapa Nuke}
\input{modely/Nuke_model.tex}
\input{matice/Nuke_matice.tex}
\input{statistiky/Nuke_stats.tex}
\newpage

\section{Mapa Overpass}
\input{modely/Overpass_model.tex}
\input{matice/Overpass_matice.tex}
\input{statistiky/Overpass_stats.tex}
\newpage

\section{Mapa Train}
\input{modely/Train_model.tex}
\input{matice/Train_matice.tex}
\input{statistiky/Train_stats.tex}
\newpage


\end{document}
